% from main.tex
%%%%%%%%%%  scan idx 271  |  folio 258  %%%%%%%%%%
tout \(\mathcal O_X\)-Module quasi-cohérent \(M\) (un tel entier existe d'après (EGA III 1.4.12 et IV 1.7.21))\,; on a alors
\[
 \parfamp_S(E)\subset[a,b]\Longrightarrow\parfamp_S(Rf_*(E))\subset[a,b+N].
\]
\end{enumerate}
Par suite, le foncteur \(Rf_*\) envoie \(D(X)_{S\text{-}\parf}\) dans \(D(Y)_{S\text{-}\parf}\)\footnote{on a d'ailleurs \(D(X)_{S\text{-}\parf}=D(X)_{S\text{-}\uparf}\) (et idem avec \(Y\)) puisque \(X\) et \(Y\) sont quasi-compacts (cf. (4.2)).}.

\emph{Preuve.} L'assertion a) découle de b), et b) s'obtient en conjuguant (2.1) et (3.7.1).

\textbf{Corollaire 4.8.1.} Soient \(f:X\to Y\) un morphisme propre et parfait, et \(E\in\Ob D^+(X)\). Alors, sous réserve de validité de (2.1) pour \(f\), si \(E\) est parfait, il en est de même de \(Rf_*(E)\)\,; en outre, si \(Y\) est quasi-compact, le foncteur \(Rf_*\) envoie \(D(X)_{\parf}\) dans \(D(Y)_{\parf}\)\footnote{On a \(D(X)_{\uparf}=D(X)_{\parf}\) (et idem avec \(Y\)).}.

\emph{Preuve.} En effet, si \(E\) est parfait, \(E\) est parfait rel. à \(f\) d'après (4.5.1 a)).

\textbf{Proposition 4.9}\footnote{Ce résultat et ses corollaires ne seront pas utilisés dans la suite du présent Séminaire, mais jouent un rôle important dans la formule de Lefschetz-Verdier pour les faisceaux cohérents, voir SGA 5 Exp.\ III.}. Soient \(S\) un schéma noethérien, \(X\) et \(Y\) des \(S\)-schémas de type fini, \(f:X\to Y\) un \(S\)-morphisme compactifiable (i.e. composé d'une immersion ouverte et d'un morphisme propre), \(E\in\Ob D(X)\), \(F\in\Ob D^+(Y)\). Si \(E\) est parfait rel. à \(f\), et \(F\) parfait rel. à \(S\), \(\RHom(E,f^!(F))\) est parfait rel. à \(S\)\footnote{Pour la définition du foncteur \(f^!\), voir [H] App.\ n\textsuperscript{o}~4.}.

\emph{Preuve.} La question étant locale sur \(X\), on peut supposer que \(f\) se factorise en \(X\xrightarrow{i}X'\xrightarrow{f'}Y\), où \(i\) est une immersion fermée et \(f'\) un morphisme lisse. Si \(E\) est parfait rel. à \(f\), \(E\) est parfait rel. à \(i\) (4.4), et si \(F\) est parfait rel. à \(S\), \(f'^!(F)\) est parfait rel. à \(S\) d'après (4.5.1 b)) puisque \(f'^!(F)\) est localement isomorphe à \(f'^*(F)\) à translation près. Comme on a \(f^!\simeq i^!f'^!\), on est donc ramené à prouver l'assertion dans le cas où \(f\) est une immersion fermée. Supposons donc \(E\) parfait rel. à \(f\), i.e. \(f_*(E)\) parfait, et \(F\) parfait rel. à \(S\). Il s'agit de montrer que \(f_*\RHom(E,f^!F)\) est
